\chapter{AgentOS 多层循环统一理论}
\chaptervoice{这一章，我们围绕“AgentOS 多层循环统一理论”做一件克制的事：把直觉压缩为状态、关系、时间尺度和可检验后果。它在全书中的任务，是说明固定功能拓扑中的持续主体。}

\section*{理论来路与依赖接口}
本章的直接思想来源是原文关于 CCM、调度、边界、卸载、单主体与 L0—L6 多层循环的 AgentOS 讨论。我们采用原文较晚形成的定义来整理较早讨论，不把对话中的试探性比喻与最终模型并列成互相竞争的“定律”。已有理论锚点主要是持续学习、工作空间、观测器—控制器与混杂系统（Kalman、混杂系统理论与 Cowan）。

本章使用上一章“单一主体原则与 Agent 进程模型”已经得到的对象与边界；本章产出将供下一章“固定拓扑持续学习的统一闭环”使用。若后文术语偶尔出现，只作为路线提示，不参与本章推导。

\section*{从哪里出发}
我们已经拥有上一章给出的概念，现在只向前走一步。我会先把对象放进闭环，再讨论它的数学形式，最后检查工程边界。读者不必预先相信本章术语；只需暂时接受一个方法论约定：智能结构的意义来自它在现实反馈中的作用，而不是来自名称本身。

令系统在观察窗口 $[t,t+T]$ 内的状态轨迹为 $x_{[t,t+T]}$，环境扰动为 $d$，行动为 $u$。本章讨论的对象若确实有效，就应改变条件分布 $p(x_{t+T}\mid x_t,u,d)$，或改变达到同一结果所需的代价。这个起点把讨论锚定在持续学习、工作空间、观测器—控制器与混杂系统上。
\conceptfigure{images/fig-l0-l6.png}{AgentOS L0—L6 多时间尺度闭环（概念图占位）}
\section{L0：现实交互闭环}
\begin{narration}
现在把镜头移到“L0：现实交互闭环”。我们只沿本章已经取得的概念向前一步：先确定它在闭环里的角色，再检查它的尺度与失败边界。
\end{narration}

本节把“L0：现实交互闭环”落实为在本章闭环中具有明确输入、输出、状态和时间尺度的功能关系。这个定义不是说“任何相似现象都算”，而是要求我们同时给出对象域、允许的干预以及观察窗口。对于主题“AgentOS 多层循环统一理论”而言，它承担的是固定功能拓扑中的持续主体中的一个可辨识环节。

把这一关系写成最小数学骨架，可得

\begin{equation}z_{t+1}=F(z_t,i_t;\alpha),\qquad o_t=H(z_t).\end{equation}

式中的符号只承担结构角色，具体参数仍需由任务和身体数据辨识。我们首先检查可观测量，再检查控制量，最后检查比它更慢的约束。这样的顺序来自持续学习、工作空间、观测器—控制器与混杂系统，但本书对L0：现实交互闭环的命名和组合仍是需要实验验证的建模提案。

这一小节的反例同样重要：若不能通过干预改变轨迹分布或资源代价，该关系在当前尺度上不可辨识。因此评价它不能只看一次任务结果，而要比较干预前后的轨迹、恢复时间、维护成本和风险。如果改变该机制而所有允许观察下的分布都不变，那么它至少在当前实验族中是冗余描述。

\begin{thought}
对“L0：现实交互闭环”做一次消融：冻结它、删除它，或只允许它以相邻时间尺度交换残差。哪些现象立即消失，哪些只在长期出现？若答案无法区分三种干预，我们还没有找到正确的状态变量。
\end{thought}

最后，把结论交还现实：本节只建立功能层关系，不由此推出特定脑区实现、主观意识或宇宙目的。可测状态、干预后果和明确失败条件，才是从原文直觉走向可检验理论的桥。

\section{L1：运行时认知控制闭环}
\begin{narration}
现在把镜头移到“L1：运行时认知控制闭环”。我们只沿本章已经取得的概念向前一步：先确定它在闭环里的角色，再检查它的尺度与失败边界。
\end{narration}

本节把“L1：运行时认知控制闭环”落实为在本章闭环中具有明确输入、输出、状态和时间尺度的功能关系。这个定义不是说“任何相似现象都算”，而是要求我们同时给出对象域、允许的干预以及观察窗口。对于主题“AgentOS 多层循环统一理论”而言，它承担的是固定功能拓扑中的持续主体中的一个可辨识环节。

把这一关系写成最小数学骨架，可得

\begin{equation}z_{t+1}=F(z_t,i_t;\alpha),\qquad o_t=H(z_t).\end{equation}

式中的符号只承担结构角色，具体参数仍需由任务和身体数据辨识。我们首先检查可观测量，再检查控制量，最后检查比它更慢的约束。这样的顺序来自持续学习、工作空间、观测器—控制器与混杂系统，但本书对L1：运行时认知控制闭环的命名和组合仍是需要实验验证的建模提案。

这一小节的反例同样重要：若不能通过干预改变轨迹分布或资源代价，该关系在当前尺度上不可辨识。因此评价它不能只看一次任务结果，而要比较干预前后的轨迹、恢复时间、维护成本和风险。如果改变该机制而所有允许观察下的分布都不变，那么它至少在当前实验族中是冗余描述。

\begin{thought}
对“L1：运行时认知控制闭环”做一次消融：冻结它、删除它，或只允许它以相邻时间尺度交换残差。哪些现象立即消失，哪些只在长期出现？若答案无法区分三种干预，我们还没有找到正确的状态变量。
\end{thought}

最后，把结论交还现实：本节只建立功能层关系，不由此推出特定脑区实现、主观意识或宇宙目的。可测状态、干预后果和明确失败条件，才是从原文直觉走向可检验理论的桥。

\section{L2：h 工作记忆认知闭环}
\begin{narration}
现在把镜头移到“L2：h 工作记忆认知闭环”。我们只沿本章已经取得的概念向前一步：先确定它在闭环里的角色，再检查它的尺度与失败边界。
\end{narration}

本节把“L2：h 工作记忆认知闭环”落实为瞬时状态、带时间索引的经验轨迹与轨迹生成结构之间的双向相变。这个定义不是说“任何相似现象都算”，而是要求我们同时给出对象域、允许的干预以及观察窗口。对于主题“AgentOS 多层循环统一理论”而言，它承担的是固定功能拓扑中的持续主体中的一个可辨识环节。

把这一关系写成最小数学骨架，可得

\begin{equation}h_t\rightarrow E_{[t-\Delta,t]}\xrightarrow{\mathsf C}\Theta.\end{equation}

式中的符号只承担结构角色，具体参数仍需由任务和身体数据辨识。我们首先检查可观测量，再检查控制量，最后检查比它更慢的约束。这样的顺序来自持续学习、工作空间、观测器—控制器与混杂系统，但本书对L2：h 工作记忆认知闭环的命名和组合仍是需要实验验证的建模提案。

这一小节的反例同样重要：把三相做成三个独立数据库会切断写入、召回和巩固的动力联系。因此评价它不能只看一次任务结果，而要比较干预前后的轨迹、恢复时间、维护成本和风险。如果改变该机制而所有允许观察下的分布都不变，那么它至少在当前实验族中是冗余描述。

\begin{thought}
对“L2：h 工作记忆认知闭环”做一次消融：冻结它、删除它，或只允许它以相邻时间尺度交换残差。哪些现象立即消失，哪些只在长期出现？若答案无法区分三种干预，我们还没有找到正确的状态变量。
\end{thought}

最后，把结论交还现实：本节只建立功能层关系，不由此推出特定脑区实现、主观意识或宇宙目的。可测状态、干预后果和明确失败条件，才是从原文直觉走向可检验理论的桥。

\section{L3：h—E—Θ 三相记忆闭环}
\begin{narration}
现在把镜头移到“L3：h—E—Θ 三相记忆闭环”。我们只沿本章已经取得的概念向前一步：先确定它在闭环里的角色，再检查它的尺度与失败边界。
\end{narration}

本节把“L3：h—E—Θ 三相记忆闭环”落实为瞬时状态、带时间索引的经验轨迹与轨迹生成结构之间的双向相变。这个定义不是说“任何相似现象都算”，而是要求我们同时给出对象域、允许的干预以及观察窗口。对于主题“AgentOS 多层循环统一理论”而言，它承担的是固定功能拓扑中的持续主体中的一个可辨识环节。

把这一关系写成最小数学骨架，可得

\begin{equation}h_t\rightarrow E_{[t-\Delta,t]}\xrightarrow{\mathsf C}\Theta.\end{equation}

式中的符号只承担结构角色，具体参数仍需由任务和身体数据辨识。我们首先检查可观测量，再检查控制量，最后检查比它更慢的约束。这样的顺序来自持续学习、工作空间、观测器—控制器与混杂系统，但本书对L3：h—E—Θ 三相记忆闭环的命名和组合仍是需要实验验证的建模提案。

这一小节的反例同样重要：把三相做成三个独立数据库会切断写入、召回和巩固的动力联系。因此评价它不能只看一次任务结果，而要比较干预前后的轨迹、恢复时间、维护成本和风险。如果改变该机制而所有允许观察下的分布都不变，那么它至少在当前实验族中是冗余描述。

\begin{thought}
对“L3：h—E—Θ 三相记忆闭环”做一次消融：冻结它、删除它，或只允许它以相邻时间尺度交换残差。哪些现象立即消失，哪些只在长期出现？若答案无法区分三种干预，我们还没有找到正确的状态变量。
\end{thought}

最后，把结论交还现实：本节只建立功能层关系，不由此推出特定脑区实现、主观意识或宇宙目的。可测状态、干预后果和明确失败条件，才是从原文直觉走向可检验理论的桥。

\section{L4：Goal—Process—Action—Verify—Closure}
\begin{narration}
现在把镜头移到“L4：Goal—Process—Action—Verify—Closure”。我们只沿本章已经取得的概念向前一步：先确定它在闭环里的角色，再检查它的尺度与失败边界。
\end{narration}

本节把“L4：Goal—Process—Action—Verify—Closure”落实为在本章闭环中具有明确输入、输出、状态和时间尺度的功能关系。这个定义不是说“任何相似现象都算”，而是要求我们同时给出对象域、允许的干预以及观察窗口。对于主题“AgentOS 多层循环统一理论”而言，它承担的是固定功能拓扑中的持续主体中的一个可辨识环节。

把这一关系写成最小数学骨架，可得

\begin{equation}z_{t+1}=F(z_t,i_t;\alpha),\qquad o_t=H(z_t).\end{equation}

式中的符号只承担结构角色，具体参数仍需由任务和身体数据辨识。我们首先检查可观测量，再检查控制量，最后检查比它更慢的约束。这样的顺序来自持续学习、工作空间、观测器—控制器与混杂系统，但本书对L4：Goal—Process—Action—Verify—Closure的命名和组合仍是需要实验验证的建模提案。

这一小节的反例同样重要：若不能通过干预改变轨迹分布或资源代价，该关系在当前尺度上不可辨识。因此评价它不能只看一次任务结果，而要比较干预前后的轨迹、恢复时间、维护成本和风险。如果改变该机制而所有允许观察下的分布都不变，那么它至少在当前实验族中是冗余描述。

\begin{thought}
对“L4：Goal—Process—Action—Verify—Closure”做一次消融：冻结它、删除它，或只允许它以相邻时间尺度交换残差。哪些现象立即消失，哪些只在长期出现？若答案无法区分三种干预，我们还没有找到正确的状态变量。
\end{thought}

最后，把结论交还现实：本节只建立功能层关系，不由此推出特定脑区实现、主观意识或宇宙目的。可测状态、干预后果和明确失败条件，才是从原文直觉走向可检验理论的桥。

\section{L5：Experience—Capability—Authority—Autonomy 生命周期闭环}
\begin{narration}
现在把镜头移到“L5：Experience—Capability—Authority—Autonomy 生命周期闭环”。我们只沿本章已经取得的概念向前一步：先确定它在闭环里的角色，再检查它的尺度与失败边界。
\end{narration}

本节把“L5：Experience—Capability—Authority—Autonomy 生命周期闭环”落实为在规定时延内具有高于普通预测控制权限的约束与中断包络。这个定义不是说“任何相似现象都算”，而是要求我们同时给出对象域、允许的干预以及观察窗口。对于主题“AgentOS 多层循环统一理论”而言，它承担的是固定功能拓扑中的持续主体中的一个可辨识环节。

把这一关系写成最小数学骨架，可得

\begin{equation}x_t\in\mathcal S,\quad u_t\in\mathcal U_{\rm safe}(x_t).\end{equation}

式中的符号只承担结构角色，具体参数仍需由任务和身体数据辨识。我们首先检查可观测量，再检查控制量，最后检查比它更慢的约束。这样的顺序来自持续学习、工作空间、观测器—控制器与混杂系统，但本书对L5：Experience—Capability—Authority—Autonomy 生命周期闭环的命名和组合仍是需要实验验证的建模提案。

这一小节的反例同样重要：让实时危险等待高层认知会违反最坏时延要求。因此评价它不能只看一次任务结果，而要比较干预前后的轨迹、恢复时间、维护成本和风险。如果改变该机制而所有允许观察下的分布都不变，那么它至少在当前实验族中是冗余描述。

\begin{thought}
对“L5：Experience—Capability—Authority—Autonomy 生命周期闭环”做一次消融：冻结它、删除它，或只允许它以相邻时间尺度交换残差。哪些现象立即消失，哪些只在长期出现？若答案无法区分三种干预，我们还没有找到正确的状态变量。
\end{thought}

最后，把结论交还现实：本节只建立功能层关系，不由此推出特定脑区实现、主观意识或宇宙目的。可测状态、干预后果和明确失败条件，才是从原文直觉走向可检验理论的桥。

\section{L6：Ψ—Goal Preference—Ω/DGA 价值与生成先验闭环}
\begin{narration}
现在把镜头移到“L6：Ψ—Goal Preference—Ω/DGA 价值与生成先验闭环”。我们只沿本章已经取得的概念向前一步：先确定它在闭环里的角色，再检查它的尺度与失败边界。
\end{narration}

本节把“L6：Ψ—Goal Preference—Ω/DGA 价值与生成先验闭环”落实为跨经历保持身份连续、同时允许受约束慢漂移的不变量集合。这个定义不是说“任何相似现象都算”，而是要求我们同时给出对象域、允许的干预以及观察窗口。对于主题“AgentOS 多层循环统一理论”而言，它承担的是固定功能拓扑中的持续主体中的一个可辨识环节。

把这一关系写成最小数学骨架，可得

\begin{equation}\dot\Psi=\varepsilon_\Psi\Pi_{\mathcal I}g(E,\Theta,\Psi).\end{equation}

式中的符号只承担结构角色，具体参数仍需由任务和身体数据辨识。我们首先检查可观测量，再检查控制量，最后检查比它更慢的约束。这样的顺序来自持续学习、工作空间、观测器—控制器与混杂系统，但本书对L6：Ψ—Goal Preference—Ω/DGA 价值与生成先验闭环的命名和组合仍是需要实验验证的建模提案。

这一小节的反例同样重要：更新过快导致身份碎裂，完全冻结则阻断必要成长。因此评价它不能只看一次任务结果，而要比较干预前后的轨迹、恢复时间、维护成本和风险。如果改变该机制而所有允许观察下的分布都不变，那么它至少在当前实验族中是冗余描述。

\begin{thought}
对“L6：Ψ—Goal Preference—Ω/DGA 价值与生成先验闭环”做一次消融：冻结它、删除它，或只允许它以相邻时间尺度交换残差。哪些现象立即消失，哪些只在长期出现？若答案无法区分三种干预，我们还没有找到正确的状态变量。
\end{thought}

最后，把结论交还现实：本节只建立功能层关系，不由此推出特定脑区实现、主观意识或宇宙目的。可测状态、干预后果和明确失败条件，才是从原文直觉走向可检验理论的桥。

\section{横向 Body/Tool/Environment/Other Agent 外化环}
\begin{narration}
现在把镜头移到“横向 Body/Tool/Environment/Other Agent 外化环”。我们只沿本章已经取得的概念向前一步：先确定它在闭环里的角色，再检查它的尺度与失败边界。
\end{narration}

本节把“横向 Body/Tool/Environment/Other Agent 外化环”落实为在本章闭环中具有明确输入、输出、状态和时间尺度的功能关系。这个定义不是说“任何相似现象都算”，而是要求我们同时给出对象域、允许的干预以及观察窗口。对于主题“AgentOS 多层循环统一理论”而言，它承担的是固定功能拓扑中的持续主体中的一个可辨识环节。

把这一关系写成最小数学骨架，可得

\begin{equation}z_{t+1}=F(z_t,i_t;\alpha),\qquad o_t=H(z_t).\end{equation}

式中的符号只承担结构角色，具体参数仍需由任务和身体数据辨识。我们首先检查可观测量，再检查控制量，最后检查比它更慢的约束。这样的顺序来自持续学习、工作空间、观测器—控制器与混杂系统，但本书对横向 Body/Tool/Environment/Other Agent 外化环的命名和组合仍是需要实验验证的建模提案。

这一小节的反例同样重要：若不能通过干预改变轨迹分布或资源代价，该关系在当前尺度上不可辨识。因此评价它不能只看一次任务结果，而要比较干预前后的轨迹、恢复时间、维护成本和风险。如果改变该机制而所有允许观察下的分布都不变，那么它至少在当前实验族中是冗余描述。

\begin{thought}
对“横向 Body/Tool/Environment/Other Agent 外化环”做一次消融：冻结它、删除它，或只允许它以相邻时间尺度交换残差。哪些现象立即消失，哪些只在长期出现？若答案无法区分三种干预，我们还没有找到正确的状态变量。
\end{thought}

最后，把结论交还现实：本节只建立功能层关系，不由此推出特定脑区实现、主观意识或宇宙目的。可测状态、干预后果和明确失败条件，才是从原文直觉走向可检验理论的桥。

\section{社会与组织超慢层}
\begin{narration}
现在把镜头移到“社会与组织超慢层”。我们只沿本章已经取得的概念向前一步：先确定它在闭环里的角色，再检查它的尺度与失败边界。
\end{narration}

本节把“社会与组织超慢层”落实为跨主体 CSO 流通过角色、协议、外部记忆和制度沉降形成的慢尺度协调结构。这个定义不是说“任何相似现象都算”，而是要求我们同时给出对象域、允许的干预以及观察窗口。对于主题“AgentOS 多层循环统一理论”而言，它承担的是固定功能拓扑中的持续主体中的一个可辨识环节。

把这一关系写成最小数学骨架，可得

\begin{equation}\mathcal O=(\{A_i\},\mathcal P,\mathcal M,\mathcal N).\end{equation}

式中的符号只承担结构角色，具体参数仍需由任务和身体数据辨识。我们首先检查可观测量，再检查控制量，最后检查比它更慢的约束。这样的顺序来自持续学习、工作空间、观测器—控制器与混杂系统，但本书对社会与组织超慢层的命名和组合仍是需要实验验证的建模提案。

这一小节的反例同样重要：把组织类比成单一大脑会虚构不存在的中央观察器。因此评价它不能只看一次任务结果，而要比较干预前后的轨迹、恢复时间、维护成本和风险。如果改变该机制而所有允许观察下的分布都不变，那么它至少在当前实验族中是冗余描述。

\begin{thought}
对“社会与组织超慢层”做一次消融：冻结它、删除它，或只允许它以相邻时间尺度交换残差。哪些现象立即消失，哪些只在长期出现？若答案无法区分三种干预，我们还没有找到正确的状态变量。
\end{thought}

最后，把结论交还现实：本节只建立功能层关系，不由此推出特定脑区实现、主观意识或宇宙目的。可测状态、干预后果和明确失败条件，才是从原文直觉走向可检验理论的桥。

\section{跨层约束与残差上浮}
\begin{narration}
现在把镜头移到“跨层约束与残差上浮”。我们只沿本章已经取得的概念向前一步：先确定它在闭环里的角色，再检查它的尺度与失败边界。
\end{narration}

本节把“跨层约束与残差上浮”落实为在本章闭环中具有明确输入、输出、状态和时间尺度的功能关系。这个定义不是说“任何相似现象都算”，而是要求我们同时给出对象域、允许的干预以及观察窗口。对于主题“AgentOS 多层循环统一理论”而言，它承担的是固定功能拓扑中的持续主体中的一个可辨识环节。

把这一关系写成最小数学骨架，可得

\begin{equation}z_{t+1}=F(z_t,i_t;\alpha),\qquad o_t=H(z_t).\end{equation}

式中的符号只承担结构角色，具体参数仍需由任务和身体数据辨识。我们首先检查可观测量，再检查控制量，最后检查比它更慢的约束。这样的顺序来自持续学习、工作空间、观测器—控制器与混杂系统，但本书对跨层约束与残差上浮的命名和组合仍是需要实验验证的建模提案。

这一小节的反例同样重要：若不能通过干预改变轨迹分布或资源代价，该关系在当前尺度上不可辨识。因此评价它不能只看一次任务结果，而要比较干预前后的轨迹、恢复时间、维护成本和风险。如果改变该机制而所有允许观察下的分布都不变，那么它至少在当前实验族中是冗余描述。

\begin{thought}
对“跨层约束与残差上浮”做一次消融：冻结它、删除它，或只允许它以相邻时间尺度交换残差。哪些现象立即消失，哪些只在长期出现？若答案无法区分三种干预，我们还没有找到正确的状态变量。
\end{thought}

最后，把结论交还现实：本节只建立功能层关系，不由此推出特定脑区实现、主观意识或宇宙目的。可测状态、干预后果和明确失败条件，才是从原文直觉走向可检验理论的桥。

\section*{本章收束：把名词还原为闭环}
现在回头看“AgentOS 多层循环统一理论”，最重要的不是记住缩写，而是说清本章对象怎样改变系统的状态、代价或可恢复性。下一章“固定拓扑持续学习的统一闭环”只使用这里已经给出的结果，不反过来为本章提供前提。

\begin{bookproposition}
若一个候选机制既不改变可观测轨迹分布，也不改变资源、风险或恢复代价，那么在当前观察尺度上，它不是一个可辨识的功能结构。
\end{bookproposition}
\begin{proofidea}
功能等价由可干预后果定义。若所有允许干预下的输出分布与代价均不变，则该机制相对于当前实验族不可辨识；要么扩大观察尺度，要么删除这一冗余描述。
\end{proofidea}

\begin{readercheck}
请尝试回答：本章对象的状态是什么？它的最快和最慢时间常数是什么？什么观测能证伪它？它失败时，残差应向哪一层升级？
\end{readercheck}
\cleardoublepage
